\section{Resonant mean motions}\label{sec:resonant}

The Diophantine condition enters the proofs of Theorems~\ref{thm:planar} and~\ref{thm:spatial} only through the
stopped drift estimate of Lemma~\ref{lem:scaling}. We replace it by Nekhoroshev's theorem for a convex Hamiltonian
with slow variables. The exponents and the mechanism are Nekhoroshev's: the Kepler Hamiltonian is convex
\cite[\S1.18]{Nek77}, and he proved the absence of collisions over exponentially long times, with the same
exponents, on his domains of planetary motions \cite[Theorem~12.3, pp.~55--56]{Nek77}. What is added is the
sharp hypothesis $C(0)\le B(\Lambda(0))-\delta$ on the whole curved sublevel.

\subsection{The input theorem}

Write points of $\R^{2(n+n_s)}$ as $(\theta,\xi;I,\eta)$, with fast angles $\theta$, fast actions $I\in\R^n$ and
slow pairs $(\xi,\eta)\in\R^{n_s}\times\R^{n_s}$, and the form $\sum dI_i\wedge d\theta_i+\sum d\eta_k\wedge
d\xi_k$. For a set $X$ and $d>0$ let $X\ominus d=\{x\in X:\{\norm{y-x}_\infty\le d\}\subset X\}$; $\aE{},\bE{}$ are
as in Theorem~\ref{thm:resonant}.

\begin{proposition}[Nekhoroshev, definite case with slow variables]\label{thm:partialnek}
Let $n\ge2$, $n_s\ge1$, let $G\subset G'\subset\R^n$ be open and convex, $S\subset\R^{2n_s}$ open and
$\rho_*>0$, with the closed $\rho_*$-ball about every point of $G$ contained in $G'$. Let $h$ be holomorphic on
$G'_{\rho_*}=\{I\in\C^n:\operatorname{dist}_\infty(I,G')<\rho_*\}$, real on real points, with
$\norm{D^2h}\le M_2$ there, and let $g,m_h>0$ and $\sigma\in\{1,-1\}$ be such that $\nrm{\nabla h(I)}\ge g$ and
$\sigma D^2h(I)(v,v)\ge m_h\abs v^2$ for all real $I\in G'_{\rho_*}$ and $v\in\R^n$. Then there exists $\mu_0>0$,
depending only on $(n,n_s,G,G',S,\rho_*,h,g,m_h,M_2)$, with the following property.

Let $0<\varepsilon<\mu_0$, and let $P$ be holomorphic on the set $F$ of
$(\theta,\xi;I,\eta)\in\C^{2(n+n_s)}$ with $\operatorname{Re}I\in G'$, $(\operatorname{Re}\xi,\operatorname{Re}\eta)
\in S$ and all imaginary parts of modulus $<\rho_*$, real on real points, $2\pi$-periodic in each $\theta_i$, with $\abs P\le1$ on $F$. Let $C>0$ and let
$(\theta,\xi,I,\eta)$ solve Hamilton's equations for $h(I)+\varepsilon P$ on $[0,C]$, with $I(t)\in G$ there and
$I(0)\in G\ominus\varepsilon^{\bE{}}$. Then $\nrm{I(t)-I(0)}\le\frac12\varepsilon^{\bE{}}$ for every
$t\in[0,\min\{C,\varepsilon^{-1}\exp(\varepsilon^{-\aE{}})\}]$ such that $(\xi(u),\eta(u))\in S\ominus\varepsilon^{\bE{}}$
for all $u\in[0,t]$.
\end{proposition}

The threshold precedes $P$, $C$ and the solution. This is \cite[Theorem~4.4, p.~30]{Nek77}, with $\sup\abs{H_1}$ as allowed by \cite[Remark~4.5.A]{Nek77}, specialized to a
definite Hessian, which is quasiconvex with all steepness indices one \cite[\S1.11, p.~9]{Nek77}; then
\cite[(4.1), (4.4), (4.5)]{Nek77} give $\zeta=n(n-1)/2$ and the exponents $\aE{},\bE{}$. It is stated in a slightly
weaker form: analyticity on an enlargement $G'$ of $G$, erosions in the maximum norm (stronger than the Euclidean
ones of \cite[\S4.2]{Nek77}), the conclusion $\le\frac12\varepsilon^{\bE{}}$ in the maximum norm, and $\mu_0$
depending on the domains. It is proved in this form in Lean (Appendix~\ref{app:formal}).

\subsection{A finite cover}

The threshold of Proposition~\ref{thm:partialnek} is not quantitative in $G'\times S$, while the slow radius of
\eqref{eq:hyp} varies with $B(\Lambda)$. We use finitely many cells, each with a slow ball at its own sharp level.
Let $n_s=n$ in the plane and $n_s=2n$ in space, and take $U$, $r_g$, $\Omega_B$ as in Section~\ref{sec:planar}
and $r_*,s,b,M$ from Lemma~\ref{lem:domain} or~\ref{lem:spdomain}, renaming the slightly smaller domain there so
that $f$ is holomorphic, real, $2\pi$-periodic in $\lambda$ and bounded by $M$ on the domain \eqref{eq:domain}
centred at any $\Lambda^0\in\Lset_\delta$. Put $h(\Lambda)=-\sum_i\nu_i^3/(2\Lambda_i^2)$ and
\[
  w=\min\{r_*,4s,4b\},\qquad\eta=w/4,\qquad\rho_*=w/8,\qquad B_0=\max\{M,1\},\qquad P=f/B_0.
\]

\begin{lemma}[Finite cover]\label{lem:cover}
There are $c_1,\dots,c_J\in\Lset_\delta$ with $\Lset_\delta\subset\bigcup_j\{\nrm{\Lambda-c_j}<\eta/2\}$ such that
the cells $G_j=\{\nrm{\Lambda-c_j}<\eta\}$, $G_j'=\{\nrm{\Lambda-c_j}<2\eta\}$ and
$S_j=\{z\in\R^{2n_s}:\frac12\abs z^2<B(c_j)-5\delta/16\}$ satisfy, with constants independent of $j$:
\begin{enumerate}[label=(\roman*)]
\item $(G_j,G_j',S_j,\rho_*,\pm h)$ satisfy the hypotheses of Proposition~\ref{thm:partialnek} with
  $\sigma=\mp1$, and $\pm P$ satisfy them on the domain $F_j$ built from $(G_j',S_j,\rho_*)$, on whose real points
  the chart Hamiltonian is $h+B_0\mu P$;
\item if $\Lambda^0\in\Lset_\delta$, $\nrm{\Lambda^0-c_j}<\eta/2$, $\Lambda\in U$ and $\nrm{\Lambda-\Lambda^0}\le4r_g$,
  then $B(\Lambda)\le B(c_j)+\delta/8$;
\item there is $\varepsilon_D>0$ with $\{\frac12\abs z^2<B(c_j)-3\delta/8\}\subset S_j\ominus r$ for all $j$ and
  $0<r<\varepsilon_D$.
\end{enumerate}
\end{lemma}

\begin{proof}
The cover exists by compactness, and $\eta+\rho_*<2\eta$. (i) Points of $(G_j')_{\rho_*}$ and of $F_j$ have
actions within $2\eta+\rho_*<w\le r_*$ of $c_j$, so with real parts in $U$; there $h$ is holomorphic and real,
$D^2h=\operatorname{diag}(-3\nu_i^3/\Lambda_i^4)$ is bounded, and at real points $\omega_i\ge g$ and
$-D^2h(v,v)\ge m_h\abs v^2$ for constants $g,m_h>0$ depending on $U$ and $w$. Points of $F_j$ also have angles
with imaginary parts below $\rho_*<4s$ and slow coordinates within $\rho_*<4b$ of
$S_j\subset\mathcal Z_0(c_j)$, so $F_j$ lies in the domain \eqref{eq:domain} centred at $c_j$, and the splitting
$H_\mu=h+\mu f$ gives the last claim. (ii) As $\eta/2\le r_g$,
$B(\Lambda)\le B(\Lambda^0)+\Omega_B(4r_g)\le B(c_j)+2\Omega_B(4r_g)\le B(c_j)+\delta/8$. (iii) With
$A_j=B(c_j)-5\delta/16$ and $A_j'=B(c_j)-3\delta/8\ge5\delta/8$: if $\frac12\abs z^2<A_j'$ and $\nrm{z'-z}\le r$,
then $\abs{z'}<\sqrt{2A_j'}+\sqrt{2n_s}\,r$; take $\varepsilon_D=\min_j(\sqrt{2A_j}-\sqrt{2A_j'})/\sqrt{2n_s}$.
\end{proof}

\subsection{Proof of Theorem~\ref{thm:resonant}}

\emph{Constants.} Proposition~\ref{thm:partialnek} gives, for each $j$, thresholds $\mu_j^\pm$ for $\pm h$. Let
$r_1>0$ be such that $nr+\Omega_B(r)\le\delta/4$ for $r\le r_1$. Choose $\varepsilon_0\in(0,1]$ with
$\varepsilon_0\le\min_j\mu_j^\pm$, $\varepsilon_0^{\bE{}}\le\min\{\eta/4,r_1,r_g\}$ and
$\varepsilon_0^{\bE{}}<\varepsilon_D$, and let $\mu_0\le\varepsilon_0/(2B_0)$ satisfy also the conditions of
Lemma~\ref{lem:osc}. For $0<\mu\le\mu_0$ put $\varepsilon=B_0\mu<\varepsilon_0$, $r=\varepsilon^{\bE{}}$ and
$T=\exp(\varepsilon^{-\aE{}})=T_0\exp(c\mu^{-\aE{}})$, with $T_0=1$, $c=B_0^{-\aE{}}$; all this precedes the
initial state.

\emph{Stopped drift.} Lemma~\ref{lem:exit} uses the normal form only through this estimate: if $\abs t\le T$ and
for all $u$ between $0$ and $t$ the state is prograde and satisfies \eqref{eq:stopped}, then
$\nrm{\Lambda(t)-\Lambda^0}\le r/2$. To prove it, orient or rotate as in Theorems~\ref{thm:planar}
and~\ref{thm:spatial} and pass to the chart, where $\frac12\abs{z}^2=C$. As $\Lambda^0\in\Lset_\delta$, choose
$j$ with $\nrm{\Lambda^0-c_j}<\eta/2$. Since $r\le\eta/4$, $\Lambda^0\in G_j\ominus r$ and $\Lambda(u)\in G_j$.
By Lemma~\ref{lem:cover}(ii),
\[
  \tfrac12\abs{z(u)}^2=C(u)<B(\Lambda(u))-\delta/2\le B(c_j)-3\delta/8,
\]
so $z(u)\in S_j\ominus r$ by Lemma~\ref{lem:cover}(iii); the bound $C<B(\Lambda^0)-\delta/4$ of
Lemma~\ref{lem:exit} would leave no room for this erosion. The orbit thus stays in the real points of $F_j$, where it
solves Hamilton's equations for $h+\varepsilon P$. For $t>0$, as $\varepsilon\le1$,
$t\le\exp(\varepsilon^{-\aE{}})\le\varepsilon^{-1}\exp(\varepsilon^{-\aE{}})$, and Proposition~\ref{thm:partialnek}
with $C=t$ gives $\nrm{\Lambda(t)-\Lambda^0}\le\frac12\varepsilon^{\bE{}}=r/2$. For $t<0$, the time-reversed
orbit on $[0,-t]$ solves Hamilton's equations for $-h-\varepsilon P$; use $\mu_j^-$.


\emph{Conclusion.} With $r$ and $T$ as above, $nr+\Omega_B(r)\le\delta/4$ and $r\le r_g$, and the improvement to
$r/2$ still gives the margin $3\delta/4$, so the proof of Lemma~\ref{lem:exit} applies unchanged. It gives existence
and uniqueness on $[-T,T]$, the margin $\delta/2$, the separations, and the drift bounds with
$A_\Lambda=B_0^{\bE{}}/2$; Lemmas~\ref{lem:osc} and~\ref{lem:normalization} give the rest. In space, $n_s=2n$ and
$h$ depends only on the $n$ fast actions, so the exponents are unchanged. Table~\ref{tab:comparison} compares the results. \qed

\begin{table}[ht]
\centering\small
\begin{tabular}{@{}lll@{}}
\toprule
 & Theorems~\ref{thm:planar}, \ref{thm:spatial} & Theorem~\ref{thm:resonant}\\
\midrule
arithmetic hypothesis & Diophantine $(\gamma,\tau)$ at $\omega(\Lambda(0))$ & none\\
constants depend on & $(n,\nu,\Aset,\delta,\gamma,\tau)$ & $(n,\nu,\Aset,\delta)$\\
time & $T_0\exp(c\mu^{-1/(2(\tau+1))})$ & $T_0\exp(c\mu^{-\aE{}})$\\
action, axis, AMD drift & $O(\sqrt\mu)$ & $O(\mu^{\bE{}})$\\
other conclusions & sharp $B$, margin $\delta/2$ & identical\\
\bottomrule
\end{tabular}
\caption{The Diophantine and the resonance-tolerant results; neither implies the other.}\label{tab:comparison}
\end{table}
